% Springer Nature render source generated from main_numerical_style.tex.
% Scientific content is preserved; only class/front-matter/build wiring is changed.
\documentclass[pdflatex,sn-mathphys-num]{sn-jnl}

\usepackage{graphicx}%
\usepackage{amsmath,amssymb,amsfonts}%
\usepackage{amsthm}%
\usepackage{mathrsfs}%
\usepackage{xcolor}%
\usepackage{textcomp}%
\usepackage{booktabs}%
\usepackage{tabularx}%
\usepackage{float}% retained to compile existing table/figure placement markup
\IfFileExists{placeins.sty}{\usepackage{placeins}}{\newcommand{\FloatBarrier}{\clearpage}}%

\floatstyle{ruled}
\newfloat{algorithm}{tbp}{loa}
\floatname{algorithm}{Algorithm}
\makeatletter
\providecommand*{\toclevel@algorithm}{1}
\newcounter{ALG@line}
\newenvironment{algorithmic}[1][0]{%
  \setcounter{ALG@line}{0}%
  \def\ALG@freq{#1}%
  \begin{list}{\ifnum\ALG@freq>0\refstepcounter{ALG@line}\arabic{ALG@line}:\fi}{%
    \setlength{\leftmargin}{2.2em}%
    \setlength{\labelwidth}{1.8em}%
    \setlength{\labelsep}{0.4em}%
    \setlength{\itemsep}{2pt}%
    \setlength{\parsep}{0pt}%
    \setlength{\topsep}{3pt}%
  }%
}{\end{list}}
\newcommand{\Require}{\item[\textbf{Require:}]}
\newcommand{\Ensure}{\item[\textbf{Ensure:}]}
\newcommand{\State}{\item}
\newcommand{\For}[1]{\State \textbf{for} #1 \textbf{do}}
\newcommand{\EndFor}{\State \textbf{end for}}
\makeatother

\theoremstyle{thmstyleone}%
\newtheorem{theorem}{Theorem}[section]
\newtheorem{proposition}[theorem]{Proposition}
\newtheorem{lemma}[theorem]{Lemma}

\theoremstyle{thmstylethree}%
\newtheorem{assumption}[theorem]{Assumption}

\theoremstyle{thmstyletwo}%
\newtheorem{remark}[theorem]{Remark}
\providecommand*{\theHassumption}{\theHtheorem}

\newcommand{\dd}{\mathrm d}
\newcommand{\E}{\mathbb E}
\newcommand{\Q}{\mathbb Q}
\newcommand{\R}{\mathbb R}
\newcommand{\F}{\mathcal F}
\providecommand{\supp}{\operatorname{supp}}
\providecommand{\footnoteA}[1]{\footnote{#1}}
\let\burl\url
\hypersetup{hypertexnames=false}

\makeatletter
\renewenvironment{table}[1][]{%
  \par\addvspace{12pt}%
  \noindent\begin{minipage}{\textwidth}%
  \def\@captype{table}%
  \centering%
  \tablebodyfont%
}{%
  \end{minipage}%
  \par\addvspace{12pt}%
}
\renewenvironment{tablenotes}[1][]{%
  \begin{list}{}{\setlength{\labelsep}{0pt}%
  \setlength{\labelwidth}{0pt}%
  \setlength{\leftmargin}{0pt}%
  \setlength{\rightmargin}{0pt}%
  \setlength{\topsep}{-6pt}%
  \setlength{\itemsep}{2pt}%
  \setlength{\partopsep}{0pt}%
  \setlength{\listparindent}{0em}%
  \setlength{\parsep}{0pt}}%
}{%
  \end{list}%
}
\makeatother

\raggedbottom

\begin{document}

\title[Hybrid LSMC--PDE for a Generalized Gatheral Double Mean-Reverting Model]{Hybrid LSMC--PDE for a Generalized Gatheral Double Mean-Reverting Model}

\author[1]{\fnm{Mara Kalicanin} \sur{Dimitrov}}\email{mara.kalicanin.dimitrov@mdu.se}
\author*[1]{\fnm{Ying} \sur{Ni}}\email{ying.ni@mdu.se}

\affil*[1]{\orgdiv{Department of Business and Mathematics}, \orgname{M\"{a}lardalen University}, \orgaddress{\postcode{721~23}, \city{V\"{a}ster{\aa}s}, \country{Sweden}}}

\abstract{
We study Bermudan put pricing under a generalized Gatheral double mean-reverting
stochastic-volatility model. The model extends the square-root specification by
allowing constant-elasticity-of-variance volatility-of-volatility coefficients while
preserving the double mean-reverting structure of the variance factors. This
generalization creates a boundary issue at zero, so we first prove strong
well-posedness and nonnegativity of the two-factor variance system for elasticity
exponents between one half and one.
We then use the model's one-way coupling structure to construct a hybrid pricing
method that combines least-squares Monte Carlo simulation with a one-dimensional
partial differential equation step. Conditional on simulated volatility-driver
information, each Bermudan continuation step reduces to a one-dimensional backward
Kolmogorov equation in log-price space. The resulting conditional asset-price step
is evaluated by a Fourier representation, and least-squares regression across the
volatility states is used to approximate the remaining continuation dependence.
We also prove the regression consistency result needed for this construction. For
the truncated regression problem used in the analysis, the fitted least-squares
coefficients converge almost surely. Consequently, for each fixed asset level, the
corresponding fitted continuation values converge uniformly over the truncated
volatility domain. Numerical experiments for Bermudan puts under a baseline
specification compare the direct hybrid estimator with plain least-squares Monte
Carlo simulation in fixed-step and fixed-path sweeps, illustrating how conditioning
on the volatility drivers affects the continuation-value regression.
}

\keywords{Bermudan option pricing, stochastic volatility, Hybrid LSMC--PDE, least-squares Monte Carlo, Gatheral double mean reversion}

\maketitle

\section{Introduction}

Bermudan option pricing under stochastic volatility combines an optimal-stopping
problem with the approximation of conditional continuation values. In deterministic
methods this usually leads to variational inequalities or parabolic
Partial Differential Equations (PDEs) on the full state space, whose cost grows
rapidly when several volatility factors are present.
Simulation methods avoid a full deterministic grid, but they shift the main difficulty
to the regression of continuation values. This difficulty is particularly visible in
stochastic-volatility models, where the continuation surface depends jointly on the
asset level and on the volatility state.

This paper studies Bermudan put pricing under a generalized Gatheral-type double
mean-reverting stochastic-volatility model. The instantaneous variance \(v\)
mean-reverts to a stochastic variance level \(v'\), while \(v'\) mean-reverts to a
long-run level. We replace the square-root volatility-of-volatility terms in the
standard double mean-reverting specification by CEV-type powers \(v^{\delta_1}\)
and \((v')^{\delta_2}\), with \(\delta_1,\delta_2\in[1/2,1]\). The square-root
model is recovered when \(\delta_1=\delta_2=1/2\), whereas larger elasticities
allow a different local volatility-of-volatility response.

The CEV generalization creates a boundary issue at zero. For exponents below one
the variance diffusion coefficients are not Lipschitz at the boundary, while the
model is meaningful only if the variance factors remain nonnegative. We therefore
first prove strong existence, pathwise uniqueness, and nonnegativity of the
two-factor variance system for \(\delta_1,\delta_2\in[1/2,1]\). The lower bound
is the natural Yamada--Watanabe threshold for power diffusion coefficients. The
proof uses the triangular structure of the variance subsystem: \(v'\) is solved
first, then \(v\) is solved conditionally on \(v'\), and the asset equation is
recovered by a stochastic exponential.

The second part of the paper uses the one-way coupling of the model to derive a
Hybrid LSMC--PDE pricing scheme. Since the variance subsystem is autonomous, we
can condition on the volatility-driver information over each Bermudan step.
After orthogonalizing the Brownian drivers, the remaining asset uncertainty is
one-dimensional, and the conditional continuation step becomes a backward
Kolmogorov equation in log-price space. This conditioning is used only inside an
iterated expectation and does not enlarge the exercise filtration.

The conditional step depends on each sampled variance segment only through a
small number of path statistics: the correlated Brownian shift, the integrated
drift, the integrated residual variance, and the terminal volatility state. We
then verify the structural assumptions of the hybrid convergence framework and
obtain almost-sure convergence of the stabilized regression coefficients for the
exact truncated regression problem. This result does not include
time-discretization, Fourier-truncation, interpolation, or finite-grid errors.

The numerical section compares the direct hybrid estimator with plain LSMC for
Bermudan puts under a baseline numerical specification.

Section~2 reviews the related literature. Section~3 introduces the model and the
Bermudan stopping problem. Section~4 proves well-posedness and nonnegativity.
Section~5 derives the conditional one-dimensional PDE representation. Section~6
verifies convergence of the truncated hybrid regression, Section~7 describes the
implementation, Section~8 reports numerical experiments, and Section~9 concludes.







\section{Related literature}

Classical diffusion pricing identifies option values with solutions of parabolic
PDEs through the Black--Scholes--Merton framework and the Feynman--Kac formula
\cite{black1973pricing,merton1973theory,karatzas1991brownian}. With early
exercise, the pricing problem becomes an optimal-stopping problem, leading to
free-boundary or variational-inequality formulations for American options
\cite{merton1977americanput,jaillet1990variational} and to a discrete Snell
recursion for Bermudan options. Tree methods such as \cite{cox1979option} and
deterministic PDE methods are effective in low dimension, but their cost
increases rapidly when several volatility or interest-rate factors are present.

Simulation methods avoid a full deterministic grid and are therefore natural in
higher-dimensional settings \cite{glasserman2004montecarlo}. For American and
Bermudan options, the main numerical difficulty is the approximation of
continuation values. The least-squares Monte Carlo method of Longstaff and
Schwartz \cite{longstaff2001valuing} and the regression-based dynamic-programming
method of Tsitsiklis and Van Roy \cite{tsitsiklis2001regression} estimate these
conditional expectations from simulated paths; convergence of least-squares
regression in this setting was analyzed in \cite{clement2002analysis}. In
stochastic-volatility models the regression problem is more demanding because the
continuation value depends jointly on the asset and volatility states.

Mixed MC--PDE methods reduce this burden by simulating selected state variables
and solving a lower-dimensional PDE for the remaining uncertainty. Such methods
for stochastic-volatility models appear in
\cite{loeper2009mixed,lipp2013mixing,cozma2017mixed}. For Bermudan options,
Farahany, Jackson, and Jaimungal \cite{farahany2020mixing} introduced a Hybrid
LSMC--PDE framework for one-way coupled stochastic-volatility models, in which
volatility factors are simulated first, the conditional asset-price uncertainty
is treated by a lower-dimensional PDE, and regression is performed across
volatility states. Further algorithmic background and multifactor examples are
discussed in \cite{farahany2019thesis}.

On the modeling side, the Heston model \cite{heston1993closedform}, based on a
Cox--Ingersoll--Ross square-root variance process \cite{cir1985termstructure},
is a standard one-factor stochastic-volatility model. Multifactor
stochastic-volatility models can improve the fit to option data
\cite{christoffersen2009smirk}. Gatheral's volatility-surface framework
\cite{gatheral2006volsurface} motivates richer volatility-factor dynamics,
including double mean-reverting specifications \cite{bayer2013fast}, where the
instantaneous variance mean-reverts to a stochastic variance level and that level
mean-reverts to a long-run variance.

The present paper keeps the double mean-reverting structure but replaces the
square-root volatility-of-volatility terms by CEV-type powers. This connects the
model to the classical CEV specification \cite{cox1975cev} and the CKLS family
of mean-reverting diffusions \cite{chan1992empirical}. The resulting variance
equations have non-Lipschitz diffusion coefficients at the origin, so
well-posedness cannot be obtained from a global-Lipschitz theorem. The
Yamada--Watanabe criterion identifies $\delta\ge 1/2$ as the natural
pathwise-uniqueness threshold for power diffusion coefficients
\cite{yamada1971uniqueness}. Unlike reflected Gatheral-type models
\cite{mishura2025skorokhod}, the present model does not add a reflection term;
instead, we prove that the nonnegative state space is preserved by the CEV-type
double mean-reverting variance system itself.

Thus the contribution is the model-specific combination of a CEV-type double
mean-reverting variance system, strong well-posedness and nonnegativity, the
conditional one-dimensional Hybrid LSMC--PDE reduction, and verification of the
truncated hybrid regression framework for Bermudan pricing.



\section{Model and Bermudan stopping problem}\label{sec:model}

We now specify the generalized gDMR model and the associated Bermudan stopping
problem. The variance factors evolve autonomously, while the asset price depends
on them; this one-way coupling is used in Section~\ref{sec:conditional_pde} to
obtain the conditional one-dimensional PDE representation.

\subsection{Stochastic basis and generalized gDMR dynamics}

Fix a finite horizon $T>0$. Let
\[
(\Omega,\F,(\F_t)_{t\in[0,T]},\Q)
\]
be a filtered probability space satisfying the usual conditions. We work under a
pricing measure $\Q$ with deterministic money-market account $B_t=e^{rt}$. The
space supports a three-dimensional correlated $(\F_t)$-Brownian motion
\[
W=(W^{(1)},W^{(2)},W^{(3)})
\]
with constant quadratic covariations
\[
\langle W^{(i)},W^{(j)}\rangle_t=\rho_{ij}t,\qquad 1\le i,j\le 3.
\]
The matrix $\rho=(\rho_{ij})_{1\le i,j\le 3}$ is symmetric positive
semidefinite and satisfies $\rho_{ii}=1$. We assume $|\rho_{23}|<1$, so that
the covariance matrix of the two volatility drivers is invertible in the
projection used in Section~\ref{sec:conditional_pde}. Because $|\rho_{23}|<1$,
the full correlation matrix is singular exactly in the degenerate case where
conditioning on the volatility drivers leaves no residual Brownian uncertainty
in the asset direction. No independence between the asset and volatility Brownian
drivers is assumed.

Let $S_0>0$, $v_0\ge0$, and $v'_0\ge0$. Motivated by Gatheral-type double mean-reverting volatility structure \citep{gatheral2006volsurface,bayer2013fast}, we consider the generalized gDMR dynamics
\begin{equation}\label{eq:gdmr}
\begin{aligned}
\dd S_t &= r\,S_t\,\dd t + S_t\sqrt{v_t}\,\dd W^{(1)}_t,\\
\dd v_t &= \kappa_1\,(v'_t - v_t)\,\dd t + \xi_1 v_t^{\delta_1}\,\dd W^{(2)}_t,\\
\dd v'_t &= \kappa_2\,(\theta - v'_t)\,\dd t + \xi_2 (v'_t)^{\delta_2}\,\dd W^{(3)}_t,
\end{aligned}
\end{equation}
where $r\ge0$, $\kappa_1,\kappa_2,\theta,\xi_1,\xi_2>0$, and
$\delta_1,\delta_2\in[1/2,1]$. The factor $v$ is the instantaneous variance and
mean-reverts to the stochastic level $v'$, while $v'$ mean-reverts to the
long-run level $\theta$. The exponents $\delta_1$ and $\delta_2$ replace the
square-root volatility-of-volatility coefficients by CEV-type powers; the
square-root double mean-reverting model is recovered when
$\delta_1=\delta_2=1/2$.

\begin{assumption}[Pricing martingale]\label{ass:pricing-martingale}
The discounted stock price $e^{-rt}S_t$ is a true $\Q$-martingale on $[0,T]$.
\end{assumption}

This condition is imposed as part of the risk-neutral pricing specification. The
well-posedness result in Section~\ref{sec:wellposed} ensures that the stochastic
exponential is well defined and that the variance factors remain nonnegative,
but it does not replace the martingale assumption.

The model itself is formulated on the nonnegative variance state space
\[
\mathcal D:=(0,\infty)\times[0,\infty)^2.
\]
For the proof of strong well-posedness in Section~\ref{sec:wellposed}, we
extend only the variance diffusion coefficients to the real line by setting
\[
\sigma_i(x):=\xi_i(x^+)^{\delta_i},\qquad
x^+:=\max\{x,0\},\qquad i=1,2.
\]
Theorem~\ref{thm:wellposed} shows that the variance factors remain
nonnegative. Hence this auxiliary extension is inactive along the solution paths
of the model in \eqref{eq:gdmr}.

\begin{remark}[One-way coupling]\label{rem:one-way}
The variance subsystem $(v,v')$ is autonomous. Hence, over one exercise
interval, the volatility-driver increments can be simulated before the residual
asset-price innovation. Conditioning on these increments fixes the realized
variance segment and the part of the asset Brownian motion correlated with the
volatility drivers; only the orthogonal asset Brownian component remains random.
This conditioning is used only inside the one-step conditional expectation and
does not enlarge the exercise filtration.
\end{remark}

\subsection{Bermudan option valuation as a discrete-time stopping problem}

Let $\pi=\{0=t_0<\dots<t_M=T\}$ be the Bermudan exercise grid, and define
\[
\mathcal T_\pi:=
\{\tau:\Omega\to\pi:\{\tau=t_n\}\in\F_{t_n},\ n=0,\dots,M\}.
\]
At the exercise date $t_n$, the payoff is $h_n(S_{t_n})$, where $h_n:(0,\infty)\to\R$ is Borel. We assume
\[
\E^\Q\!\left[\max_{0\le n\le M} e^{-rt_n}|h_n(S_{t_n})|\right]<\infty.
\]
This condition is immediate for bounded payoffs, including the Bermudan put
payoff $h_n(s)=(K-s)^+$. No compact-support condition is imposed at this stage;
compactly supported truncations enter only in the regression convergence
argument of Section~\ref{sec:as_convergence}.

The well-posedness result in Section~\ref{sec:wellposed} implies uniqueness in
law for the time-homogeneous Markov process $(S_t,v_t,v'_t)$ on $\mathcal D$.
Hence the Bermudan value functions may be chosen as Borel functions. Write
\(\Delta t_n:=t_{n+1}-t_n\), and set
\[
    U_M(s,v,w):=h_M(s),
\]
where \(w\) denotes the second variance coordinate. For \(n=M-1,\ldots,0\),
define a Borel version of the continuation function by
\[
C_n(s,v,w)
:=e^{-r\Delta t_n}\,\E^\Q\!\left[
U_{n+1}\bigl(S_{t_{n+1}},v_{t_{n+1}},v'_{t_{n+1}}\bigr)
\middle|\,
S_{t_n}=s,\,v_{t_n}=v,\,v'_{t_n}=w
\right],
\]
and set
\[
U_n(s,v,w):=\max\{h_n(s),C_n(s,v,w)\}.
\]
The adapted process
\[
Y_n:=U_n(S_{t_n},v_{t_n},v'_{t_n}),\qquad n=0,\dots,M,
\]
then satisfies the current-value Snell recursion
\begin{equation}\label{eq:snell}
\begin{aligned}
Y_M&=h_M(S_{t_M}),\\
Y_n&=\max\Bigl\{h_n(S_{t_n}),\,
e^{-r\Delta t_n}\,
\E^\Q\!\bigl[Y_{n+1}\mid\F_{t_n}\bigr]\Bigr\},
\qquad n=M-1,\dots,0.
\end{aligned}
\end{equation}
Thus $e^{-rt_n}Y_n$ is the Snell envelope of the discounted payoff process on
the grid. The Bermudan time-zero value is
\[
V_0
=\sup_{\tau\in\mathcal T_\pi}\E^\Q\!\left[e^{-r\tau}h_\tau(S_\tau)\right]
=U_0(S_0,v_0,v'_0),
\]
where $h_\tau(S_\tau)$ denotes $h_n(S_{t_n})$ on the event $\{\tau=t_n\}$. The continuation term in \eqref{eq:snell} is therefore
\[
e^{-r\Delta t_n}\,\E^\Q\!\bigl[Y_{n+1}\mid\F_{t_n}\bigr]
=C_n\bigl(S_{t_n},v_{t_n},v'_{t_n}\bigr)\qquad\text{$\Q$-a.s.}
\]




\section{Well-posedness of the generalized variance system}\label{sec:wellposed}

Section~\ref{sec:model} formulates the generalized gDMR dynamics on the nonnegative
variance state space and uses the corresponding transition kernel in the
Bermudan recursion. We now verify stochastic invariance of this state space. The only
non-Lipschitz coefficients are the CEV-type variance diffusions. The triangular
structure of the model allows us to solve $v'$ first, then $v$ conditional on
$v'$, and finally the asset equation.

\begin{assumption}[CEV exponents]\label{ass:delta-range}
Throughout this section we impose
\[
\delta_1,\delta_2\in[1/2,1].
\]
\end{assumption}

We shall use the following elementary consequence of the Yamada--Watanabe
criterion. For $\delta\in[1/2,1]$, the map $x\mapsto (x^+)^\delta$ satisfies
\[
\bigl|(x^+)^\delta-(y^+)^\delta\bigr|
\le |x-y|^\delta,\qquad x,y\in\R.
\]
Hence the diffusion coefficient $x\mapsto \xi(x^+)^\delta$ satisfies the
Yamada--Watanabe modulus condition, because the choice $\rho(u)=C u^\delta$
gives
\[
\int_0^\varepsilon \frac{\dd u}{\rho(u)^2}
=\infty
\qquad\text{for }\delta\ge \frac12.
\]
The upper bound $\delta\le1$ gives at most linear growth.

\begin{theorem}[Strong well-posedness and nonnegativity]\label{thm:wellposed}
Assume Assumption~\ref{ass:delta-range}. For every $T>0$ and every initial
state
\[
S_0>0,\qquad v_0\ge0,\qquad v'_0\ge0,
\]
the nonnegative-state system \eqref{eq:gdmr} admits a unique strong solution
$(S_t,v_t,v'_t)_{t\in[0,T]}$ with continuous paths such that
\[
S_t>0,\qquad v_t\ge0,\qquad v'_t\ge0,\qquad
0\le t\le T,\quad \Q\text{-a.s.}
\]
Moreover, for every deterministic initial state in $\mathcal D$, the solution
is a time-homogeneous strong Markov process on $\mathcal D$.
\end{theorem}

\begin{proof}
We first solve the auxiliary variance equations on $\R$, obtained by replacing
the power coefficients by their positive-part extensions. Once nonnegativity has
been proved, the auxiliary equations coincide with the original system on the
state space $\mathcal D$.

\emph{Step 1: construction of $v'$.}
Define
\[
b_2(x):=\kappa_2(\theta-x),\qquad
\sigma_2(x):=\xi_2(x^+)^{\delta_2},\qquad x\in\R.
\]
The drift $b_2$ is globally Lipschitz and has linear growth. The diffusion
coefficient $\sigma_2$ is continuous, has at most linear growth, and satisfies
the Yamada--Watanabe modulus condition by Assumption~\ref{ass:delta-range}.
Therefore pathwise uniqueness holds for
\[
\dd v'_t=b_2(v'_t)\dd t+\sigma_2(v'_t)\dd W^{(3)}_t,
\qquad v'_0\ge0;
\]
see \cite{yamada1971uniqueness} and \cite[Ch.~5]{karatzas1991brownian}. Since
the coefficients are continuous with linear growth, weak existence holds. Weak
existence together with pathwise uniqueness yields a unique strong solution by
the Yamada--Watanabe theorem.

It remains to show that the solution does not leave the nonnegative half-line.
At the boundary,
\[
\sigma_2(0)=0,\qquad b_2(0)=\kappa_2\theta\ge0.
\]
Moreover, on the negative extension one has
\[
\sigma_2(x)=0,\qquad b_2(x)=\kappa_2(\theta-x)\ge0,\qquad x\le0.
\]
Thus the diffusion has no outward noise at the boundary and the drift points
inward on the negative side. The one-dimensional stochastic invariance
criterion for closed intervals, or equivalently Tanaka's formula applied to the
negative part with this positive-part extension, gives
\[
v'_t\ge0,\qquad 0\le t\le T,\quad \Q\text{-a.s.}
\]

\emph{Step 2: construction of $v$.}
Having constructed $v'$, define
\[
b_1(t,x):=\kappa_1(v'_t-x),\qquad
\sigma_1(x):=\xi_1(x^+)^{\delta_1}.
\]
For $m\ge1$, set
\[
\tau_m:=\inf\{t\ge0:v'_t>m\}\wedge T.
\]
Since $v'$ has continuous paths on $[0,T]$, we have $\tau_m\uparrow T$ almost
surely. On $[0,\tau_m]$, the drift $b_1$ is progressively measurable,
Lipschitz in $x$ with constant $\kappa_1$, and satisfies
\[
|b_1(t,x)|\le\kappa_1(m+|x|).
\]
The diffusion coefficient $\sigma_1$ is continuous, has at most linear growth,
and satisfies the Yamada--Watanabe modulus condition. The localized
one-dimensional existence and pathwise-uniqueness theorem for SDEs with
progressive random coefficients, obtained by the same Yamada--Watanabe argument
after localization, therefore gives a unique strong stopped solution of
\[
\dd v_t=b_1(t,v_t)\dd t+\sigma_1(v_t)\dd W^{(2)}_t,
\qquad t\le\tau_m.
\]
The possible correlation between $W^{(2)}$ and $W^{(3)}$ is immaterial here:
pathwise uniqueness compares solutions driven by the same Brownian vector and
the same realized input path $v'$. By pathwise uniqueness, the stopped solutions
are compatible on overlaps. Letting $m\to\infty$ gives a unique strong solution
for $v$ on $[0,T]$.

The same invariance argument gives nonnegativity. Since $v'_t\ge0$,
\[
b_1(t,0)=\kappa_1v'_t\ge0,\qquad \sigma_1(0)=0.
\]
On the negative extension,
\[
b_1(t,x)=\kappa_1(v'_t-x)\ge0,\qquad
\sigma_1(x)=0,\qquad x\le0.
\]
Hence the closed half-line is invariant and
\[
v_t\ge0,\qquad 0\le t\le T,\quad \Q\text{-a.s.}
\]

\emph{Step 3: construction of $S$.}
The process $v$ is continuous and finite on the compact interval $[0,T]$.
Therefore
\[
\int_0^T v_s\,\dd s<\infty\qquad \Q\text{-a.s.}
\]
The asset equation is linear in $S$, and its unique strong solution is
\[
S_t
=S_0
\exp\!\left(
\int_0^t \Bigl(r-\tfrac12 v_s\Bigr)\dd s
+\int_0^t \sqrt{v_s}\,\dd W_s^{(1)}
\right).
\]
Thus $S_t>0$ for all $t\in[0,T]$ whenever $S_0>0$.

The auxiliary positive-part extension is inactive along the constructed variance
paths, because $v_t,v'_t\ge0$. Hence the constructed solution solves the
original nonnegative-state system \eqref{eq:gdmr}. Pathwise uniqueness for the
auxiliary triangular system gives pathwise uniqueness for \eqref{eq:gdmr}.

Finally, existence and pathwise uniqueness for every deterministic initial state
imply uniqueness in law. Since the coefficients are time homogeneous, the
standard strong Markov property for well-posed time-homogeneous SDEs gives that
the solution is a time-homogeneous strong Markov process on $\mathcal D$.
\end{proof}

Consequently, the state process and the transition kernel used in the Bermudan
recursion are well defined. The martingale property of the discounted stock
remains the pricing assumption in Assumption~\ref{ass:pricing-martingale}. The
condition $\delta_i\ge1/2$ is the Yamada--Watanabe pathwise-uniqueness threshold
for power diffusion coefficients; no Feller-type condition is imposed, because
the theorem proves nonnegativity rather than strict positivity or boundary
inaccessibility.

\section{Conditional one-dimensional PDE representation}\label{sec:conditional_pde}

We now derive the one-step conditional representation used by the hybrid method.
On a Bermudan interval, the autonomous variance subsystem allows us to condition
on the volatility-driver information. After projecting the asset Brownian driver
onto the volatility drivers, only the orthogonal asset Brownian component remains
random. The conditional continuation step is therefore a one-dimensional
backward Kolmogorov equation in log-price space.

\subsection{Orthogonalization on a Bermudan step}

Fix an interval $[t_n,t_{n+1}]$ in the Bermudan exercise grid and define
\[
\mathcal H_n
:=
\F_{t_n}
\vee
\sigma\!\left(
W_s^{(2)}-W_{t_n}^{(2)},
W_s^{(3)}-W_{t_n}^{(3)}
: s\in[t_n,t_{n+1}]
\right).
\]
Since the variance subsystem is autonomous and driven only by
$(W^{(2)},W^{(3)})$, the path
\[
(v_s,v'_s)_{s\in[t_n,t_{n+1}]}
\]
is $\mathcal H_n$-measurable. The same conditioning also fixes the part of the
asset Brownian motion that is correlated with the volatility drivers. We now
identify the remaining residual Brownian innovation.

Assume $|\rho_{23}|<1$ and set
\[
\Sigma_{23}:=
\begin{bmatrix}
1 & \rho_{23}\\
\rho_{23} & 1
\end{bmatrix},
\qquad
\beta:=\Sigma_{23}^{-1}
\begin{bmatrix}
\rho_{12}\\
\rho_{13}
\end{bmatrix}
=
\begin{bmatrix}
\beta_2\\
\beta_3
\end{bmatrix}.
\]
Equivalently,
\begin{equation}\label{eq:proj-coeffs}
\beta_2=\frac{\rho_{12}-\rho_{13}\rho_{23}}{1-\rho_{23}^2},
\qquad
\beta_3=\frac{\rho_{13}-\rho_{12}\rho_{23}}{1-\rho_{23}^2}.
\end{equation}
The residual variance in the Gaussian projection of $W^{(1)}$ onto the vector
$(W^{(2)},W^{(3)})$ is
\begin{equation}\label{eq:sigperp}
\begin{aligned}
\sigma_\perp^2
&:=
1-
\begin{bmatrix}\rho_{12}&\rho_{13}\end{bmatrix}
\Sigma_{23}^{-1}
\begin{bmatrix}\rho_{12}\\ \rho_{13}\end{bmatrix} \\
&=
\frac{
1-\rho_{12}^2-\rho_{13}^2-\rho_{23}^2
+2\rho_{12}\rho_{13}\rho_{23}
}{
1-\rho_{23}^2
}
\ge0.
\end{aligned}
\end{equation}
The nonnegativity is the Schur-complement condition implied by positive
semidefiniteness of the full correlation matrix.

When $\sigma_\perp>0$, the process
\[
B_t^\perp
:=
\frac{W_t^{(1)}-\beta_2W_t^{(2)}-\beta_3W_t^{(3)}}{\sigma_\perp}
\]
is a Brownian motion. Its increments over $[t_n,t_{n+1}]$ are independent of
$\mathcal H_n$. Hence, on $[t_n,t_{n+1}]$,
\begin{equation}\label{eq:orth}
\dd W_t^{(1)}
=\beta_2\,\dd W_t^{(2)}
+\beta_3\,\dd W_t^{(3)}
+\sigma_\perp\,\dd B_t^\perp.
\end{equation}
In the degenerate case $\sigma_\perp=0$, the last term is absent and the asset
Brownian driver is fully determined by the two volatility Brownian drivers.

Let $X_t:=\log S_t$. Then
\begin{equation}\label{eq:logS}
\dd X_t
=\Bigl(r-\tfrac12v_t\Bigr)\dd t+\sqrt{v_t}\,\dd W_t^{(1)}.
\end{equation}
Substituting the orthogonal decomposition gives
\[
\dd X_t
=\Bigl(r-\tfrac12v_t\Bigr)\dd t
+\sqrt{v_t}\Bigl(\beta_2\,\dd W_t^{(2)}
+\beta_3\,\dd W_t^{(3)}\Bigr)
+\sigma_\perp\sqrt{v_t}\,\dd B_t^\perp.
\]

Define the $\mathcal H_n$-measurable correlated shift
\[
Z_t^{\mathrm{sh}}
:=
\int_{t_n}^t
\sqrt{v_s}
\Bigl(\beta_2\,\dd W_s^{(2)}+\beta_3\,\dd W_s^{(3)}\Bigr),
\qquad t\in[t_n,t_{n+1}],
\]
and put
\[
Y_t:=X_t-Z_t^{\mathrm{sh}}.
\]
Then $X_t=Y_t+Z_t^{\mathrm{sh}}$, and, conditional on $\mathcal H_n$,
\[
\dd Y_t
=\Bigl(r-\tfrac12v_t\Bigr)\dd t
+\sigma_\perp\sqrt{v_t}\,\dd B_t^\perp.
\]
Thus, after conditioning on $\mathcal H_n$, the only remaining randomness in
the asset equation is the one-dimensional residual Brownian motion $B^\perp$.
In the degenerate case $\sigma_\perp=0$, even this residual randomness vanishes.

\subsection{Conditional expectation as a pathwise PDE}

The next result is the analytic object used by the hybrid method. For a bounded
next-step value function written in the log-price variable, it propagates the
value over one sampled variance segment. If $F$ is an asset-space value
function, the corresponding log-space terminal function is
\[
G(y,v,w):=F(e^y,v,w).
\]
The representation in this section is undiscounted; the factor
$e^{-r\Delta t_n}$ in the Bermudan recursion is applied after the conditional
PDE step.

\begin{lemma}[Conditional backward Kolmogorov equation]\label{lem:cond-PDE}
Under Assumption~\ref{ass:delta-range} and the correlation assumptions in
Section~\ref{sec:model}, let $G:\R\times[0,\infty)^2\to\R$ be bounded and Borel
measurable. Fix $n\in\{0,\dots,M-1\}$, and work on $[t_n,t_{n+1}]$ with
$\mathcal H_n$ as above. Set
\[
Z_n:=Z_{t_{n+1}}^{\mathrm{sh}}
=
\int_{t_n}^{t_{n+1}}
\sqrt{v_s}
\Bigl(\beta_2\,\dd W_s^{(2)}+\beta_3\,\dd W_s^{(3)}\Bigr).
\]
For $t\in[t_n,t_{n+1}]$ and $y\in\R$, let $Y^{t,y}$ solve
\[
\dd Y_s^{t,y}
=\Bigl(r-\tfrac12v_s\Bigr)\dd s
+\sigma_\perp\sqrt{v_s}\,\dd B_s^\perp,
\qquad
Y_t^{t,y}=y,
\qquad s\in[t,t_{n+1}].
\]
Define
\[
u(t,y)
:=
\E^\Q\!\left[
G\!\left(
Y_{t_{n+1}}^{t,y}+Z_n,
v_{t_{n+1}},
v'_{t_{n+1}}
\right)
\,\middle|\,
\mathcal H_n
\right].
\]
Then $u$ is an $\mathcal H_n$-measurable random field. Conditionally on
$\mathcal H_n$, it is the bounded mild solution of the possibly degenerate
backward Kolmogorov problem
\[
\begin{cases}
\partial_t u(t,y)
+\Bigl(r-\tfrac12v_t\Bigr)\partial_y u(t,y)
+\frac12\sigma_\perp^2v_t\,\partial_{yy}u(t,y)
=0,
&
(t,y)\in[t_n,t_{n+1})\times\R,
\\[0.4em]
u(t_{n+1},y)
=G\!\left(y+Z_n,v_{t_{n+1}},v'_{t_{n+1}}\right),
&
y\in\R.
\end{cases}
\]
If $G$ is bounded and continuous, this mild solution is also the bounded
viscosity solution. When $\sigma_\perp=0$, the equation reduces to the
corresponding first-order transport equation.
\end{lemma}

\begin{proof}
The variables $Z_n$, $v_{t_{n+1}}$, $v'_{t_{n+1}}$, and the path
$s\mapsto v_s$ on $[t_n,t_{n+1}]$ are $\mathcal H_n$-measurable. Conditional on
$\mathcal H_n$, $Y^{t,y}$ is a one-dimensional time-inhomogeneous diffusion
with deterministic coefficients. The conditional Feynman--Kac formula
\cite[Ch.~5]{karatzas1991brownian} gives the displayed mild Kolmogorov
representation. If $\sigma_\perp=0$, the second-order term is absent and the
equation reduces to deterministic transport.
\end{proof}

The previous lemma can be made more explicit because the coefficients depend on
time but not on $y$. For $t\in[t_n,t_{n+1}]$, define
\[
I_n(t)
:=
\int_t^{t_{n+1}}
\Bigl(r-\tfrac12v_s\Bigr)\dd s,
\qquad
B_n(t)
:=
\sigma_\perp^2\int_t^{t_{n+1}}v_s\,\dd s.
\]
In particular,
\[
I_n:=I_n(t_n),\qquad B_n:=B_n(t_n).
\]

\begin{proposition}[Gaussian form of the conditional step]\label{prop:conditional-gaussian-step}
Let $\xi\sim N(0,1)$ be an auxiliary variable independent of $\mathcal H_n$.
Then, conditionally on $\mathcal H_n$, the solution in
Lemma~\ref{lem:cond-PDE} satisfies
\[
u(t,y)
=
\E_\xi\!\left[
G\!\left(
y+Z_n+I_n(t)+\sqrt{B_n(t)}\,\xi,
v_{t_{n+1}},
v'_{t_{n+1}}
\right)
\right],
\]
with the convention that the Gaussian term is absent when $B_n(t)=0$.
Equivalently, when $B_n(t)>0$,
\[
u(t,y)
=
\int_{\R}
G\!\left(
y+Z_n+I_n(t)+\sqrt{B_n(t)}\,z,
v_{t_{n+1}},
v'_{t_{n+1}}
\right)
\phi_N(z)\,\dd z,
\]
where
\[
\phi_N(z):=(2\pi)^{-1/2}e^{-z^2/2}.
\]
When $B_n(t)=0$,
\[
u(t,y)
=
G\!\left(y+Z_n+I_n(t),v_{t_{n+1}},v'_{t_{n+1}}\right).
\]
\end{proposition}

\begin{proof}
Conditional on $\mathcal H_n$, the solution of the residual equation is
\[
Y_{t_{n+1}}^{t,y}
=
y+I_n(t)
+\int_t^{t_{n+1}}\sigma_\perp\sqrt{v_s}\,\dd B_s^\perp.
\]
The stochastic integral is Gaussian with mean zero and variance $B_n(t)$, since
the integrand is deterministic after conditioning and the residual Brownian
increments are independent of $\mathcal H_n$. This gives the stated formula.
\end{proof}

Thus, for a next-step asset-space value function $F$, the pathwise pre-surface
is
\[
\bar C_n(s;\mathcal H_n)
:=
e^{-r\Delta t_n}u(t_n,\log s).
\]
The true continuation value is obtained by averaging $\bar C_n(s;\mathcal H_n)$
over the volatility-driver information conditional on the current state. Hence
each sampled segment enters the asset-direction propagation only through
\[
Z_n,\qquad I_n,\qquad B_n,\qquad v_{t_{n+1}},\qquad v'_{t_{n+1}},
\]
with $(v_{t_n},v'_{t_n})$ used in the outer regression.


\section{Convergence of the truncated hybrid regression}
\label{sec:as_convergence}

Section~\ref{sec:conditional_pde} shows that, on each Bermudan interval,
conditioning on the volatility-driver information produces a pathwise
conditional continuation value in the asset direction. The remaining step is to
approximate the dependence of this pathwise value on the current volatility
state by least-squares regression. This section verifies, for the generalized
double mean-reverting model, the structural assumptions required by the Hybrid
LSMC--PDE convergence theorem of \cite{farahany2020mixing}.

The result below concerns the exact truncated regression problem. It assumes
that the conditional one-step expectations are evaluated exactly through the
pathwise PDE representation. It does not include time-discretization error in
the variance simulation, Fourier truncation error, interpolation error, or
finite-grid error in the asset coordinate.

Throughout this section, write
\[
    V_n:=(v_{t_n},v'_{t_n}),
    \qquad
    \Delta t_n:=t_{n+1}-t_n .
\]

\subsection{Regression set-up and truncation conditions}
\label{subsec:regression_truncation}

Let $d_B\in\mathbb{N}$ and let
\[
    \phi = (\phi_1,\dots,\phi_{d_B})
    : [0,\infty)^2 \to \R^{d_B}
\]
be a fixed vector of basis functions in the two-factor volatility state $(v,v')$.
Let $\mathcal D_v\subset [0,\infty)^2$ be a compact rectangle with
$\supp(\phi)\subset \mathcal D_v$. Let
$\mathcal S_{\mathrm{tr}}\subset(0,\infty)$ be a compact asset truncation
domain. All payoff and continuation functions in this section are understood as
bounded continuous compactly supported truncations on $(0,\infty)$, equal to the
intended functions on the numerical asset region and extended outside that
region by the chosen cutoff.

Let $\Q'$ denote the product probability law under which each cross-sectional
sample is an i.i.d.\ copy of the joint per-path object under $\Q$: the
volatility-driver segment, its exercise-date states, and the associated path
statistics used in the regression. For any bounded Borel next-step function
\[
    F:(0,\infty)\times[0,\infty)^2\to\R,
\]
define the pathwise hybrid one-step operator, obtained by restarting the asset
coordinate at the input value $s$ while conditioning on the sampled
volatility-driver information, by
\[
    \bar{\mathcal C}_n^{F}(s;\mathcal H_n)
    :=
    e^{-r\Delta t_n}
    \E^{\Q}\!\left[
        F(S_{t_{n+1}},v_{t_{n+1}},v'_{t_{n+1}})
        \,\middle|\,
        S_{t_n}=s,\mathcal H_n
    \right].
\]
Here $\mathcal H_n$ is the volatility-driver information on
$[t_n,t_{n+1}]$ defined in Section~\ref{sec:conditional_pde}. Thus
$\bar{\mathcal C}_n^{F}(s;\mathcal H_n)$ is the quantity obtained after
integrating out the residual asset Brownian uncertainty, but before averaging
over the volatility-driver segment. The true continuation value associated with
$F$ is
\[
    \E^{\Q'}\!\left[
        \bar{\mathcal C}_n^{F}(s;\mathcal H_n)
        \,\middle|\,
        V_n
    \right].
\]

We define the idealized truncated value functions by backward induction. Set
\[
    f_M(s,v,v'):=h_M(s).
\]
For $n=M-1,\ldots,1$, suppose $f_{n+1}$ is already defined. For fixed
$s\in\mathcal S_{\mathrm{tr}}$, set
\[
\Gamma_{n,s}
:=
\bar{\mathcal C}_n^{f_{n+1}}(s;\mathcal H_n).
\]
The population least-squares objective is
\begin{equation}
\begin{aligned}
H_{n,s}(a)
&:=
\E^{\Q'}\!\left[
    \left(
        \Gamma_{n,s}
        -
        a\cdot\phi(V_n)
    \right)^2
\right],
\qquad a\in\R^{d_B}.
\end{aligned}
\label{eq:Hns}
\end{equation}
Equivalently, since $\phi(V_n)$ is $V_n$-measurable, the same normal equations
are obtained by replacing $\Gamma_{n,s}$ with
$\E^{\Q'}[\Gamma_{n,s}\mid V_n]$. Thus the regression may be viewed either as a
projection of the pathwise hybrid target or as a projection of the true
continuation value onto the span of the basis functions.

Define
\begin{equation}
\begin{aligned}
A_n
&:=
\E^{\Q'}\!\left[
    \phi(V_n)\phi(V_n)^\top
\right],\\
b_n(s)
&:=
\E^{\Q'}\!\left[
    \phi(V_n)\Gamma_{n,s}
\right].
\end{aligned}
\label{eq:ideal_coeff_normal_eq}
\end{equation}
When $A_n$ is nonsingular, the idealized coefficient vector is
\[
    a_n(s)=A_n^{-1}b_n(s),
\]
and the corresponding truncated continuation and value functions are
\[
\begin{aligned}
    c_n(s,v,v')&:=a_n(s)\cdot\phi(v,v'),\\
    f_n(s,v,v')&:=\max\{h_n(s),c_n(s,v,v')\}.
\end{aligned}
\]

Under $\Q'$, let
\[
    \{(\mathcal H_n^{\,j},V_n^{\,j},V_{n+1}^{\,j}):j\ge1\}
\]
be the i.i.d.\ cross-sectional samples of the volatility-driver information and
the associated exercise-date volatility states at step $n$. The corresponding
statistics $Z_n^{\,j}$, $I_n^{\,j}$, and $B_n^{\,j}$ are measurable functions of
$\mathcal H_n^{\,j}$ and are evaluated along the same sampled segment. Given the
already constructed sample value function $f_{n+1}^N$, define
\[
\bar C_{n}^{\,j,N}(s)
:=
    \bar{\mathcal C}_{n,j}^{f_{n+1}^N}(s),
\]
where the subscript $j$ means that the conditional PDE step is evaluated along
the $j$-th sampled variance-driver segment. Then
\[
\begin{aligned}
    A_n^N
    &:=
    \frac1N\sum_{j=1}^N
        \phi(V_n^{\,j})\phi(V_n^{\,j})^\top,\\
    b_n^N(s)
    &:=
    \frac1N\sum_{j=1}^N
        \phi(V_n^{\,j})\bar C_n^{\,j,N}(s).
\end{aligned}
\]
For $R>0$, define the truncated inverse by
\begin{equation}
\bigl[A\bigr]^{-1}_R
:=
\begin{cases}
\bigl[A\bigr]^{-1},
&\text{if \(A\) is invertible and }\|A^{-1}\|\le R,\\
0,
&\text{otherwise.}
\end{cases}
\label{eq:estimated_coeff}
\end{equation}
The stabilized sample coefficient is
\[
    a_n^N(s):=\bigl[A_n^N\bigr]^{-1}_R b_n^N(s).
\]
Finally,
\[
\begin{aligned}
    c_n^N(s,v,v')&:=a_n^N(s)\cdot\phi(v,v'),\\
    f_n^N(s,v,v')&:=\max\{h_n(s),c_n^N(s,v,v')\},
\end{aligned}
\]
with terminal condition $f_M^N=f_M=h_M$.


\begin{assumption}[Truncation conditions]
\label{ass:truncation}
The following conditions are imposed for the convergence statement.

\begin{enumerate}
    \item The basis functions $\phi_1,\ldots,\phi_{d_B}$ are continuous,
    bounded, and compactly supported in the compact rectangle
    $\mathcal D_v\subset[0,\infty)^2$.

    \item For each exercise date $t_n$, the truncated payoff $h_n$ is
    continuous, bounded, and compactly supported in
    $\mathcal S_{\mathrm{tr}}\subset(0,\infty)$.

    \item The population Gram matrix $A_n$ is nonsingular for each
    $n=1,\ldots,M-1$.

    \item The truncation constant $R$ is chosen so that
    \[
        R>\max_{1\le n\le M-1}\|A_n^{-1}\|.
    \]
\end{enumerate}
\end{assumption}

\begin{remark}[Truncation and compact support]\label{rem:truncation-compact-support}
The compact-support assumptions belong to the truncated regression problem, not
to the original stochastic-volatility model. For a put payoff on $(0,\infty)$,
one uses a continuous cutoff near $s=0$ and, if needed, beyond a large upper
asset level, while keeping the payoff unchanged on the numerical asset region.
Similarly, basis functions that are naturally polynomial may be multiplied by a
smooth cutoff outside $\mathcal D_v$. These truncations are used to invoke the
convergence theorem; they are not part of the financial dynamics.
\end{remark}

\subsection{Verification of the hybrid structure}
\label{subsec:separability}

The Hybrid LSMC--PDE convergence theorem of \cite{farahany2020mixing} requires
multiplicative separability of the asset dynamics and finite-dimensional
dependence of the conditional one-step operator on each sampled volatility-driver
segment. Both properties hold here.

\begin{proposition}[Structural verification]
\label{prop:hybrid-assumptions}
Under Assumption~\ref{ass:delta-range}, the generalized double mean-reverting
model satisfies the positivity, one-way coupling, multiplicative separability,
and finite-dimensional path-statistic assumptions required by the Hybrid
LSMC--PDE convergence framework. More precisely, for $0\le t<u\le T$,
\[
S_u=S_tR_{t,u},
\qquad
R_{t,u}:=
\exp\left(
\int_t^u\left(r-\frac12v_s\right)\,\dd s
+
\int_t^u\sqrt{v_s}\,\dd W_s^{(1)}
\right),
\]
so $R_{t,u}$ does not depend on the value of $S_t$. Moreover, for each exercise
step $n$, the pathwise hybrid operator satisfies
\[
\bar{\mathcal C}_n^F(s;\mathcal H_n)
=
e^{-r\Delta t_n}
\E_\xi\!\left[
F\left(
s\exp(Z_n+I_n+\sqrt{B_n}\xi),
V_{n+1}
\right)
\right],
\qquad \xi\sim N(0,1),
\]
with the deterministic interpretation when $B_n=0$. Hence the sampled
volatility-driver segment enters the conditional operator only through
\[
\Theta_n:=(V_n,V_{n+1},Z_n,I_n,B_n).
\]
Under Assumption~\ref{ass:truncation}, the compactness and boundedness
assumptions of the truncated regression problem are also satisfied.
\end{proposition}

\begin{proof}
Theorem~\ref{thm:wellposed} gives well-posedness, nonnegativity of the variance
factors, and positivity of the asset process. The variance subsystem is
autonomous, so the model is one-way coupled. The stochastic exponential formula
gives the multiplicative representation of $S_u$, and
Proposition~\ref{prop:conditional-gaussian-step} gives the displayed
finite-dimensional conditional representation. The remaining compactness and
boundedness conditions are exactly those imposed in Assumption~\ref{ass:truncation}.
\end{proof}



\subsection{Main convergence theorem}
\label{subsec:main_convergence}

\begin{theorem}[Almost-sure convergence of the hybrid regression coefficients]
\label{thm:as_coeff_convergence}
Assume the generalized double mean-reverting model of
Section~\ref{sec:model} under its standing correlation assumptions, and assume
Assumption~\ref{ass:delta-range} and Assumption~\ref{ass:truncation}. Work with
the exact truncated hybrid regression problem defined above, so that the
one-step targets are the pathwise quantities
$\Gamma_{n,s}=\bar{\mathcal C}_n^{f_{n+1}}(s;\mathcal H_n)$, evaluated without
time-discretization, FFT, interpolation, or asset-grid error. Fix
\[
    n\in\{1,\ldots,M-1\}
    \qquad\text{and}\qquad
    s\in\mathcal S_{\mathrm{tr}}.
\]
Let $a_n(s)$ be the idealized coefficient vector defined by the population
normal equations, and let $a_n^N(s)$ be the stabilized hybrid estimator based on
$N$ i.i.d.\ cross-sectional samples under $\Q'$.
Then
\[
    \lim_{N\to\infty}
    \|a_n^N(s)-a_n(s)\|
    =
    0,
    \qquad
    \Q'\text{-a.s.}
\]
Consequently,
\[
    \lim_{N\to\infty}
    \sup_{(v,v')\in\mathcal D_v}
    |c_n^N(s,v,v')-c_n(s,v,v')|
    =
    0,
    \qquad
    \Q'\text{-a.s.}
\]
For this fixed asset level $s$, the corresponding truncated value functions
also satisfy
\[
    \lim_{N\to\infty}
    \sup_{(v,v')\in\mathcal D_v}
    |f_n^N(s,v,v')-f_n(s,v,v')|
    =
    0,
    \qquad
    \Q'\text{-a.s.}
\]
\end{theorem}

\begin{proof}
By Proposition~\ref{prop:hybrid-assumptions}, the model satisfies the
positivity, one-way coupling, multiplicative separability, and
finite-dimensional path-dependence assumptions of
\cite[Theorem~1]{farahany2020mixing}. Assumption~\ref{ass:truncation} gives the
boundedness, compact support, nonsingularity, and stabilized-inverse truncation
conditions. The strict choice of $R$ ensures that the population inverse lies
inside the truncation region and that the sample inverse is eventually
untruncated on the event $A_n^N\to A_n$. Therefore
\cite[Theorem~1]{farahany2020mixing} yields
\[
    \|a_n^N(s)-a_n(s)\|\to0,
    \qquad
    \Q'\text{-a.s.}
\]

The uniform convergence of $c_n^N(s,\cdot,\cdot)$ on $\mathcal D_v$ follows
from the boundedness of the basis vector on $\mathcal D_v$:
\[
\begin{aligned}
    \sup_{(v,v')\in\mathcal D_v}
    |c_n^N(s,v,v')-c_n(s,v,v')|
    &=
    \sup_{(v,v')\in\mathcal D_v}
    |(a_n^N(s)-a_n(s))\cdot\phi(v,v')| \\
    &\le
    \|a_n^N(s)-a_n(s)\|
    \sup_{(v,v')\in\mathcal D_v}\|\phi(v,v')\|.
\end{aligned}
\]
The convergence of $f_n^N$ follows from the elementary inequality
\[
    |\max\{h_n(s),x\}-\max\{h_n(s),y\}|\le |x-y|.
\]
\end{proof}

Since the numerical implementation in Section~\ref{sec:hybrid-method} uses
only finitely many asset grid points, the fixed-$s$ theorem applies
simultaneously on that grid by intersecting finitely many probability-one
events. No uniform convergence in the continuous asset variable
$s\in\mathcal S_{\mathrm{tr}}$ is claimed here.




\section{Implementation of the Hybrid LSMC--PDE scheme}
\label{sec:hybrid-method}

We now describe the numerical implementation of the Hybrid LSMC--PDE scheme for
the generalized gDMR model. The exact conditional representation was derived in
Section~\ref{sec:conditional_pde}, and the convergence of the truncated
regression problem was verified in Section~\ref{sec:as_convergence}. The present
section adds the numerical ingredients needed in practice: simulation of the
variance subsystem, accumulation of the path statistics $Z_n,I_n,B_n$, FFT
evaluation of the one-dimensional conditional propagation, and stabilized
regression over the volatility states.

The discretization choices below are part of the implementation. They are
distinct from the exact truncated regression problem analyzed in
Section~\ref{sec:as_convergence}; in particular, Euler discretization error, FFT
truncation error, interpolation error, and finite-grid effects are not included
in the convergence theorem.

\subsection{Approximation spaces}
Let
\[
    \pi=\{0=t_0<t_1<\cdots<t_M=T\}
\]
be the Bermudan exercise grid, and let $h_n$ denote the exercise payoff at
$t_n$. We work on a uniform log-price grid
\[
    y_i=y_{\min}+(i-1)\Delta y,
    \qquad
    \Delta y=\frac{y_{\max}-y_{\min}}{N_S-1},
    \qquad
    i=1,\ldots,N_S,
\]
and set
\[
    s_i=e^{y_i},
    \qquad
    \mathcal S_h:=\{s_1,\ldots,s_{N_S}\}.
\]
The interval $[y_{\min},y_{\max}]$ is chosen large enough that the relevant
asset values remain away from the artificial boundaries. In the FFT step we use
an enlarged log-price interval, or equivalently padding, to reduce wrap-around
error from the periodic convolution.

For the volatility variables, write $w=v'$. Choose a compact rectangle
\[
    \mathcal D_v=[0,v_{\max}]\times[0,w_{\max}]\subset[0,\infty)^2
\]
and a vector of continuous bounded basis functions
\[
    \phi=(\phi_1,\ldots,\phi_{d_B}),
    \qquad
    \supp(\phi_k)\subset \mathcal D_v.
\]
For numerical work, the bounds $v_{\max}$ and $w_{\max}$ may be chosen from a
preliminary simulation or from conservative parameter-based quantiles. Once chosen, the domain
$\mathcal D_v$ and the basis functions are fixed before the training regression;
data-dependent truncation selection is not part of the convergence theorem in
Section~\ref{sec:as_convergence}.

Following the separable approximation used in the hybrid regression framework,
the continuation value at $t_n$ is represented as
\[
    c_n(s,v,w)\approx a_n(s)\cdot\phi(v,w),
    \qquad
    a_n(s)\in\R^{d_B}.
\]
Numerically, the coefficient vectors are computed only at the grid points
$s_i\in\mathcal S_h$. Off-grid asset values are evaluated by interpolation in the
log-price coordinate.

\subsection{Variance simulation and accumulated statistics}
Let
\[
    0=\tau_0<\tau_1<\cdots<\tau_L=T
\]
be an Euler grid containing all exercise dates. Write
\[
    \Delta\tau_\ell:=\tau_{\ell+1}-\tau_\ell.
\]
For each path $j$, simulate correlated Gaussian increments
\[
    (\Delta W_{\ell}^{(2),j},\Delta W_{\ell}^{(3),j})
\]
with variances $\Delta\tau_\ell$ and covariance
$\rho_{23}\Delta\tau_\ell$. The numerical implementation uses the following
positive-part Euler discretization:
\[
\begin{aligned}
    \bar w_\ell^j&:=(w_\ell^j)^+,
    \qquad
    \bar v_\ell^j:=(v_\ell^j)^+,\\
    w_{\ell+1}^j
    &:=
    \left[
        w_\ell^j+\kappa_2(\theta-\bar w_\ell^j)\Delta\tau_\ell
        +\xi_2(\bar w_\ell^j)^{\delta_2}\Delta W_{\ell}^{(3),j}
    \right]^+,\\
    v_{\ell+1}^j
    &:=
    \left[
        v_\ell^j+\kappa_1(\bar w_\ell^j-\bar v_\ell^j)\Delta\tau_\ell
        +\xi_1(\bar v_\ell^j)^{\delta_1}\Delta W_{\ell}^{(2),j}
    \right]^+.
\end{aligned}
\]
Other stable nonnegative discretizations may be substituted in different
implementations, provided that the corresponding statistic accumulation is
changed consistently.

For each exercise interval $[t_n,t_{n+1}]$, the path statistics are accumulated
over the Euler substeps contained in that interval:
\[
\begin{aligned}
    Z_n^j
    &:=
    \sum_{\tau_\ell\in[t_n,t_{n+1})}
        \sqrt{\bar v_\ell^j}
        \left(
            \beta_2\Delta W_{\ell}^{(2),j}
            +
            \beta_3\Delta W_{\ell}^{(3),j}
        \right),\\
    I_n^j
    &:=
    \sum_{\tau_\ell\in[t_n,t_{n+1})}
        \left(r-\frac12\bar v_\ell^j\right)\Delta\tau_\ell,\\
    B_n^j
    &:=
    \sigma_\perp^2
    \sum_{\tau_\ell\in[t_n,t_{n+1})}
        \bar v_\ell^j\Delta\tau_\ell .
\end{aligned}
\]
Here $\beta_2,\beta_3$ are the projection coefficients in
\eqref{eq:proj-coeffs}, and $\sigma_\perp$ is given by \eqref{eq:sigperp}.
The left-point rule is natural for $Z_n^j$, since the stochastic integral is an
It\^{o} integral. The same left-point values are used for $I_n^j$ and $B_n^j$ for
consistency with the Euler simulation, and the same increments
$\Delta W_{\ell}^{(2),j},\Delta W_{\ell}^{(3),j}$ are used in the variance
updates and in $Z_n^j$.

\subsection{Pathwise FFT propagation}
Fix $n\in\{0,\ldots,M-1\}$. Suppose that $\widehat V_{n+1}^N$ has already been
constructed on the asset grid and is represented through the volatility basis.
For each Monte Carlo path $j$, the simulated volatility-driver information on
$[t_n,t_{n+1}]$ determines the endpoint state
\[
    (v_{t_{n+1}}^j,w_{t_{n+1}}^j)
\]
and the statistics $Z_n^j,I_n^j,B_n^j$. Define the terminal grid function
\[
    g_{n+1}^j(y_i)
    :=
    \widehat V_{n+1}^N
    \left(e^{y_i},v_{t_{n+1}}^j,w_{t_{n+1}}^j\right),
    \qquad i=1,\ldots,N_S.
\]
By the Gaussian representation in Section~\ref{sec:conditional_pde},
\[
    u_n^j(y)
    =
    \E_\xi\!\left[
        g_{n+1}^j
        \left(
            y+Z_n^j+I_n^j+\sqrt{B_n^j}\xi
        \right)
    \right],
    \qquad \xi\sim N(0,1).
\]
With the Fourier convention
\[
    \mathcal F f(\omega)
    =
    \int_{\R}e^{-\mathrm i\omega y}f(y)\,\dd y,
\]
we obtain
\[
    \mathcal F u_n^j(\omega)
    =
    \mathcal F g_{n+1}^j(\omega)
    \exp\left(
        \mathrm i\omega Z_n^j+\mathrm i\omega I_n^j-\frac12\omega^2B_n^j
    \right).
\]
Define
\[
    \Psi_n^j(\omega)
    :=
    \mathrm i\omega Z_n^j+\mathrm i\omega I_n^j-\frac12\omega^2B_n^j.
\]
On the uniform log-price grid, the FFT implementation is
\begin{equation}
    u_n^j(\cdot)
    \approx
    \operatorname{FFT}^{-1}
    \left(
        \operatorname{FFT}[g_{n+1}^j]\,
        \exp(\Psi_n^j)
    \right),
\label{eq:fft-recursion}
\end{equation}
where the exponential is applied pointwise at the discrete Fourier frequencies.
The FFT input is the unshifted terminal grid $g_{n+1}^j$; the correlated shift
$Z_n^j$ is included only through the multiplier. Equivalently, one may shift the
terminal grid first and omit $\mathrm i\omega Z_n^j$, but the two conventions
must not be combined. If $B_n^j=0$, the conditional propagation reduces to the
deterministic log-price shift $Z_n^j+I_n^j$.

The pathwise pre-surface on the asset grid is
\[
    \widehat C_n^j(s_i)
    :=
    e^{-r\Delta t_n}u_n^j(y_i),
    \qquad i=1,\ldots,N_S.
\]

\subsection{Regression and backward recursion}
For each exercise date $t_n$, $n=M-1,\ldots,1$, and each asset grid point $s_i$,
the regression estimates the dependence of the pathwise pre-surface
$\widehat C_n^j(s_i)$ on the current volatility state. Define
\[
    V_n^j:=(v_{t_n}^j,w_{t_n}^j).
\]
The sample Gram matrix is
\[
    A_n^N
    :=
    \frac1N\sum_{j=1}^N
        \phi(V_n^j)\phi(V_n^j)^\top,
\]
and the right-hand side for the grid point $s_i$ is
\[
    b_n^N(s_i)
    :=
    \frac1N\sum_{j=1}^N
        \phi(V_n^j)\widehat C_n^j(s_i).
\]
Using the truncated inverse from Section~\ref{sec:as_convergence}, set
\begin{equation}
    a_n^N(s_i):=[A_n^N]^{-1}_R b_n^N(s_i).
\label{eq:lsq}
\end{equation}
The fitted continuation surface is
\[
    \widehat C_n^N(s_i,v,w)
    :=
    a_n^N(s_i)\cdot\phi(v,w),
\]
and the Bermudan update is
\[
    \widehat V_n^N(s_i,v,w)
    :=
    \max\left\{
        h_n(s_i),
        \widehat C_n^N(s_i,v,w)
    \right\}.
\]
Starting from
\[
    \widehat V_M^N(s_i,v,w):=h_M(s_i),
\]
the recursion is performed backward for $n=M-1,\ldots,1$. The fitted functions
are defined at the asset grid points and extended to off-grid asset values by
log-price interpolation.

\subsection{Time-zero direct estimator}
At $t_0=0$, the initial volatility state is fixed, so no cross-sectional
regression is needed. After $\widehat V_1^N$ has been fitted, the time-zero
continuation value is evaluated by averaging first-step pathwise pre-surfaces:
\[
    \widehat C_0^{N_0}(s_i)
    :=
    \frac1{N_0}
    \sum_{j=1}^{N_0}
        \widehat C_0^{j,\mathrm{eval}}(s_i).
\]
The evaluation paths may be independent of the training paths, or the same paths
may be reused with $N_0=N$, in which case the result is an in-sample direct
estimator.
The direct estimator is
\[
    \widehat V_{d,0}(S_0)
    =
    \max\left\{
        h_0(S_0),
        \operatorname{Interp}_y\!\left[
            \widehat C_0^{N_0}
        \right](\log S_0)
    \right\}.
\]

\begin{algorithm}[H]
\caption{Hybrid LSMC--PDE recursion}\label{alg:hybrid-lsmc-pde}
\begin{algorithmic}[1]
\Require Exercise grid $\{t_n\}_{n=0}^M$, Euler grid containing the exercise dates, log-price grid $\{y_i\}_{i=1}^{N_S}$, volatility domain $\mathcal D_v$, basis $\phi$, truncation level $R$, training size $N$, and evaluation size $N_0$
\Ensure Direct estimator $\widehat V_{d,0}(S_0)$ and fitted continuation surfaces $\widehat C_n^N$
\State Simulate $N$ training paths of the autonomous variance subsystem on the Euler grid.
\State Store exercise-date volatility states and accumulate $Z_n^j$, $I_n^j$, and $B_n^j$ over each interval $[t_n,t_{n+1}]$.
\State Set $\widehat V_M^N(s_i,v,w)\gets h_M(s_i)$ on $\mathcal S_h\times\mathcal D_v$.
\For{$n=M-1,\ldots,1$}
    \For{$j=1,\ldots,N$}
        \State Compute $\{\widehat C_n^j(s_i)\}_{i=1}^{N_S}$ using the unshifted terminal grid $g_{n+1}^j$ in \eqref{eq:fft-recursion}, with $Z_n^j$ included only through the multiplier $\exp(\Psi_n^j)$; do not pre-shift $g_{n+1}^j$ by $Z_n^j$.
    \EndFor
    \For{each $s_i\in\mathcal S_h$}
        \State Compute $b_n^N(s_i)$ and $a_n^N(s_i)$ by the truncated-inverse regression \eqref{eq:lsq}.
    \EndFor
    \State Set $\widehat C_n^N(s_i,v,w)\gets a_n^N(s_i)\cdot\phi(v,w)$ for all $s_i\in\mathcal S_h$.
    \State Set $\widehat V_n^N(s_i,v,w)\gets\max\{h_n(s_i),\widehat C_n^N(s_i,v,w)\}$ for all $s_i\in\mathcal S_h$.
\EndFor
\State Form the time-zero evaluation sample: either simulate $N_0$ independent paths on $[t_0,t_1]$, or reuse the $N$ training paths with $N_0=N$ for the in-sample estimator; compute $\widehat C_0^{j,\mathrm{eval}}(s_i)$.
\State Set $\widehat C_0^{N_0}(s_i)\gets N_0^{-1}\sum_{j=1}^{N_0}\widehat C_0^{j,\mathrm{eval}}(s_i)$.
\State \textbf{return} $\widehat V_{d,0}(S_0)=\max\{h_0(S_0),\operatorname{Interp}_y[\widehat C_0^{N_0}](\log S_0)\}$.
\end{algorithmic}
\end{algorithm}
\section{Numerical experiments}\label{sec:numerical}

This section reports direct Bermudan put pricing experiments under a baseline
generalized gDMR numerical specification. The plain LSMC estimator
\cite{longstaff2001valuing} is compared with the Hybrid LSMC--PDE estimator
described in Section~7. We report direct prices and reference-relative errors.
Standard errors and confidence intervals are used for the reference table and
for the error bars in the figures.

\subsection{Experimental setting and reference values}

Table~\ref{tab:baseline-setup} reports the fixed gDMR parameter set and
numerical settings used in the experiments. Since \(S_0=100\), the strikes \(K=70\), \(K=80\), and
\(K=90\) are out-of-the-money, \(K=100\) is at-the-money, and \(K=110\) is
in-the-money for the Bermudan put.
\begin{table}[htbp]
\centering
\caption{Fixed gDMR parameter set and numerical settings used in the Bermudan pricing experiments.}
\label{tab:baseline-setup}
\begin{threeparttable}
\small
\begin{tabularx}{0.94\textwidth}{@{}l c X@{}}
\toprule
Symbol & Value & Interpretation \\
\midrule
\multicolumn{3}{@{}l}{\textit{Model and contract parameters}} \\
\(S_0\) & 100 & initial asset level \\
\(v_0\) & 0.114 & initial instantaneous variance factor \\
\(v'_0\) & 0.110 & initial secondary variance factor \\
\(r\) & 0.02 & risk-free rate \\
\(T\) & 1 & maturity in years \\
\(\kappa_1\) & 5.5 & mean-reversion speed of \(v_t\) toward \(v'_t\) \\
\(\kappa_2\) & 0.1 & mean-reversion speed of \(v'_t\) toward \(\theta\) \\
\(\theta\) & 0.078 & long-run variance level \\
\(\xi_1\) & 2.689 & volatility scale of the first variance factor \\
\(\xi_2\) & 0.502 & volatility scale of the second variance factor \\
\(\delta_1\) & 0.94 & CEV exponent of the first variance factor \\
\(\delta_2\) & 0.94 & CEV exponent of the second variance factor \\
\(\rho_{12}\) & -0.982 & correlation between the first and second Brownian drivers \\
\(\rho_{13}\) & -0.727 & correlation between the first and third Brownian drivers \\
\(\rho_{23}\) & 0.590 & correlation between the second and third Brownian drivers \\
\addlinespace
\multicolumn{3}{@{}l}{\textit{Fixed numerical setting}} \\
\(N_{\mathrm{ex}}\) & 12 & positive Bermudan exercise dates, including maturity (excluding \(t_0\)) \\
\(K\) & \(\{70,80,90,100,110\}\) & strike set used in all experiments \\
\(N_S\) & 301 & asset-grid points in the conditional PDE solver \\
\((a_{\min},a_{\max})\) & \((0.30,3.50)\) & asset-range factors relative to \(S_0\) \\
\(q_v\) & 0.999 & volatility truncation quantile \\
\((M_{\mathrm{ref}},N_{\mathrm{ref}})\) & \((1200,1{,}200{,}000)\) & reference LSMC step count and path budget \\
\botrule
\end{tabularx}
\begin{tablenotes}[flushleft]
\footnotesize
\item High-resolution reference prices are reported separately in Table~\ref{tab:baseline-reference-prices}. The Hybrid LSMC--PDE grid and truncation settings are kept fixed throughout the two study families.
\end{tablenotes}
\end{threeparttable}
\end{table}

The correlation assumptions used in Section~5 are satisfied. In particular,
\(|\rho_{23}|<1\), and for this parameter set the residual variance after
conditioning on the two volatility drivers is small but positive:
\[
    \sigma_\perp^2 \approx 2.25\times 10^{-3}.
\]
Thus the conditional PDE step still integrates out a nonzero residual asset
Brownian component.

The plain LSMC regressions use a cubic \(16\)-term basis
\[
1,x,y,z,x^2,y^2,z^2,xy,xz,yz,x^3,y^3,z^3,p,p^2,p^3,
\]
where \(x=S/S_0\), \(y=v/\theta\), \(z=v'/\theta\), and
\(p=(K-S)^+/K\). The hybrid regressions use a volatility-only cubic basis
\[
1,y,z,y^2,yz,z^2,y^3,y^2z,yz^2,z^3,
\]
where \(y\) and \(z\) are the normalized volatility coordinates after applying
the numerical truncation to the chosen volatility rectangle \(D_v\). Thus the
reported implementation uses clipped volatility coordinates in the regression;
the smooth compact-support cutoff in Section~6 belongs to the idealized
convergence setting. For the hybrid estimator, the same simulated
volatility-driver paths are used for fitting and for the time-zero direct
average; the reported hybrid direct estimator is therefore an in-sample direct
estimator in the sense described in Section~7.

For a direct estimator \(V\) and reference value \(V^{\mathrm{ref}}\), we report
\[
    \varepsilon_{\mathrm{dir}}
    :=
    \frac{|V-V^{\mathrm{ref}}|}{|V^{\mathrm{ref}}|}.
\]
The reference values are high-resolution direct LSMC estimates with \(1200\)
Euler steps and \(1{,}200{,}000\) direct paths. They are used as numerical
reference values, not as exact prices. The reported standard errors are empirical
Monte Carlo standard errors of the corresponding direct estimators, and the
95\% confidence intervals are computed as the direct estimate plus or minus
\(1.96\) standard errors. For the Hybrid LSMC--PDE estimator, these standard
errors are computed from the pathwise first-step direct values generated by the
implemented backward recursion. They should therefore be interpreted as
direct-estimator sampling errors for the implemented scheme, not as a full
numerical error decomposition including reference uncertainty, time
discretization, Fourier truncation, interpolation, or regression-training
variability. In the relative-error plots, the reference value
\(V^{\mathrm{ref}}\) is treated as fixed. The error bars are obtained by mapping
the direct-estimator confidence intervals through
\[
    V\mapsto \frac{|V-V^{\mathrm{ref}}|}{|V^{\mathrm{ref}}|}.
\]
\begin{table}[htbp]
\centering
\caption{High-resolution direct LSMC reference values for the baseline numerical specification.}
\label{tab:baseline-reference-prices}
\small
\begin{tabularx}{0.82\textwidth}{@{}cccc@{}}
\toprule
\(K\) & Reference price & SE & 95\% confidence interval \\
\midrule
70 & 2.522851 & 0.007103 & [2.508929, 2.536773] \\
80 & 4.298933 & 0.009345 & [4.280617, 4.317249] \\
90 & 6.942119 & 0.011774 & [6.919042, 6.965196] \\
100 & 10.682037 & 0.014105 & [10.654391, 10.709683] \\
110 & 15.743520 & 0.015757 & [15.712636, 15.774404] \\
\botrule
\end{tabularx}
\end{table}

\subsection{Fixed path budget and varying Euler steps}

The first experiment keeps the direct path budget fixed and varies the number of
Euler time steps. The two path budgets are \(20{,}000\) and \(60{,}000\), and
both methods are evaluated at \(24,48,72,\) and \(96\) Euler steps. This
experiment checks how the direct Bermudan estimator behaves as the variance
discretization is refined away from the high-resolution reference regime.

For \(20{,}000\) paths, the LSMC reference-error range over the full step-sweep
grid is \(0.344\%\)--\(8.645\%\), while the Hybrid LSMC--PDE range is
\(0.019\%\)--\(7.459\%\). Figure~\ref{fig:step-sweep-20k} reports the strike-wise
reference-error profiles, and Table~\ref{tab:step-sweep-20k-step72} reports
representative values at \(72\) Euler steps.

\begin{figure}[htbp]
\centering
\includegraphics[width=\textwidth]{Figures/step_sweep_20k_direct_relative_error.pdf}
\caption{Reference-relative errors for the matched \(20{,}000\)-path step-sweep experiment. Error bars are induced by the 95\% confidence intervals of the direct estimators.}
\label{fig:step-sweep-20k}
\end{figure}

\begin{table}[htbp]
\centering
\caption{Representative direct prices and reference errors at \(20{,}000\) paths and 72 Euler steps.}
\label{tab:step-sweep-20k-step72}
\small
\begin{tabular}{lrrrr}
\toprule
Strike & LSMC price & LSMC ref. err. & Hybrid price & Hybrid ref. err. \\
\midrule
\(70\) & 2.578706 & 2.214\% & 2.547003 & 0.957\% \\
\(80\) & 4.408717 & 2.554\% & 4.307120 & 0.190\% \\
\(90\) & 7.068072 & 1.814\% & 6.905102 & 0.533\% \\
\(100\) & 10.816495 & 1.259\% & 10.606747 & 0.705\% \\
\(110\) & 15.893853 & 0.955\% & 15.676443 & 0.426\% \\
\botrule
\end{tabular}
\end{table}

For \(60{,}000\) paths, the same design gives an LSMC reference-error range of
\(0.102\%\)--\(7.720\%\) and a Hybrid LSMC--PDE range of
\(0.028\%\)--\(7.857\%\). Figure~\ref{fig:step-sweep-60k} and
Table~\ref{tab:step-sweep-60k-step72} give the corresponding graphical and
tabular summaries. The hybrid estimator is not uniformly better at every coarse
or fine discretization point, but it gives smaller reference errors in most of
the reported configurations, with the clearest differences at the intermediate
step levels.

\begin{figure}[htbp]
\centering
\includegraphics[width=\textwidth]{Figures/step_sweep_60k_direct_relative_error.pdf}
\caption{Reference-relative errors for the matched \(60{,}000\)-path step-sweep experiment. Error bars are induced by the 95\% confidence intervals of the direct estimators.}
\label{fig:step-sweep-60k}
\end{figure}

\begin{table}[htbp]
\centering
\caption{Representative direct prices and reference errors at \(60{,}000\) paths and 72 Euler steps.}
\label{tab:step-sweep-60k-step72}
\small
\begin{tabular}{lrrrr}
\toprule
Strike & LSMC price & LSMC ref. err. & Hybrid price & Hybrid ref. err. \\
\midrule
\(70\) & 2.578570 & 2.209\% & 2.556746 & 1.344\% \\
\(80\) & 4.361900 & 1.465\% & 4.326225 & 0.635\% \\
\(90\) & 7.032832 & 1.307\% & 6.940184 & 0.028\% \\
\(100\) & 10.793695 & 1.045\% & 10.662148 & 0.186\% \\
\(110\) & 15.853083 & 0.696\% & 15.749100 & 0.035\% \\
\botrule
\end{tabular}
\end{table}

\FloatBarrier

\subsection{Fixed Euler step count and varying path budget}

The second experiment fixes the number of Euler time steps and varies the direct
path budget. The reported path counts are \(250,1000,5000,10000,20000,40000\),
and \(60000\). Two discretization levels, \(48\) and \(60\) Euler steps, are
considered. This experiment is aimed at the sampling and regression behavior of
the direct estimators.

At \(48\) Euler steps, the at-the-money strike \(K=100\) gives an LSMC
reference-error range of \(1.105\%\)--\(25.948\%\) over the path grid, while the
Hybrid LSMC--PDE range is \(0.088\%\)--\(11.118\%\).
Figure~\ref{fig:path-sweep-48} reports the full strike-wise path sweep, and
Table~\ref{tab:path-sweep-48-path20k} reports the representative
\(20{,}000\)-path values.

\begin{figure}[htbp]
\centering
\includegraphics[width=0.93\textwidth]{Figures/path_sweep_steps48_direct_relative_error.pdf}
\caption{Reference-relative errors for the fixed \(48\)-step path-sweep experiment. Error bars are induced by the 95\% confidence intervals of the direct estimators.}
\label{fig:path-sweep-48}
\end{figure}

\begin{table}[htbp]
\centering
\caption{Representative direct prices and reference errors at \(20{,}000\) paths and fixed 48 Euler steps.}
\label{tab:path-sweep-48-path20k}
\small
\begin{tabular}{lrrrr}
\toprule
Strike & LSMC price & LSMC ref. err. & Hybrid price & Hybrid ref. err. \\
\midrule
\(70\) & 2.611646 & 3.520\% & 2.574765 & 2.058\% \\
\(80\) & 4.432813 & 3.114\% & 4.345363 & 1.080\% \\
\(90\) & 7.129239 & 2.695\% & 6.957089 & 0.216\% \\
\(100\) & 10.806834 & 1.168\% & 10.672650 & 0.088\% \\
\(110\) & 15.920337 & 1.123\% & 15.746530 & 0.019\% \\
\botrule
\end{tabular}
\end{table}

At \(60\) Euler steps, the corresponding at-the-money LSMC range is
\(1.324\%\)--\(28.932\%\), while the Hybrid LSMC--PDE range is
\(0.049\%\)--\(10.102\%\). Figure~\ref{fig:path-sweep-60} and
Table~\ref{tab:path-sweep-60-path20k} give the companion summaries. The
reduction in reference error is most visible in the low- and moderate-path
regimes, where the plain regression estimator is most affected by sampling
noise.

\begin{figure}[htbp]
\centering
\includegraphics[width=0.93\textwidth]{Figures/path_sweep_steps60_direct_relative_error.pdf}
\caption{Reference-relative errors for the fixed \(60\)-step path-sweep experiment. Error bars are induced by the 95\% confidence intervals of the direct estimators.}
\label{fig:path-sweep-60}
\end{figure}

\begin{table}[htbp]
\centering
\caption{Representative direct prices and reference errors at \(20{,}000\) paths and fixed 60 Euler steps.}
\label{tab:path-sweep-60-path20k}
\small
\begin{tabular}{lrrrr}
\toprule
Strike & LSMC price & LSMC ref. err. & Hybrid price & Hybrid ref. err. \\
\midrule
\(70\) & 2.635093 & 4.449\% & 2.563087 & 1.595\% \\
\(80\) & 4.483550 & 4.294\% & 4.317574 & 0.434\% \\
\(90\) & 7.119832 & 2.560\% & 6.910022 & 0.462\% \\
\(100\) & 10.909701 & 2.131\% & 10.609574 & 0.678\% \\
\(110\) & 15.917591 & 1.106\% & 15.679793 & 0.405\% \\
\botrule
\end{tabular}
\end{table}

Overall, the baseline numerical experiments are consistent with the conditional
PDE step acting as a variance-reduction and dimension-reduction device for the
continuation approximation. The results are direct-estimator comparisons against
high-resolution LSMC references, not exact pricing errors. The full direct price
grids are reported in Appendix~\ref{app:direct-prices}.

\section{Conclusion}\label{sec:conclusion}
This paper studied Bermudan option pricing under a generalized Gatheral double mean-reverting stochastic volatility model, where the volatility-of-volatility coefficients are allowed to follow CEV-type power functions. This generalization extends the classical square-root specification and gives the model greater flexibility, but it also creates a non-Lipschitz boundary behavior at zero that must be addressed before the model can be used in a rigorous pricing setting.

The paper makes two main contributions. First, it establishes strong well-posedness and nonnegativity of the two-factor variance system in the exponent regime $\delta_1, \delta_2 \in [1/2,1]$. This regime is natural from the perspective of pathwise uniqueness and includes the square-root case as a special specification. These results provide the mathematical foundation required to interpret the generalized dynamics as a valid stochastic-volatility model and to formulate the Bermudan stopping problem in a consistent way.

Second, the paper shows that the generalized model fits naturally into the Hybrid LSMC--PDE framework. Since the variance factors evolve autonomously, the one-way coupling structure is preserved. Conditional on the volatility-driver information, including the correlated Brownian shift, the log-price process reduces to a one-dimensional diffusion with deterministic time-dependent coefficients, so the continuation value over each Bermudan step can be computed through a one-dimensional backward parabolic PDE. This leads to a hybrid procedure that combines pathwise PDE solves with stabilized regression across the volatility states. Within this framework, the fitted regression coefficients converge almost surely for each fixed regression date and asset level, with corresponding fixed-asset continuation convergence on the volatility truncation set.

The numerical experiments were carried out for a fixed gDMR parameter set with
CEV exponents \(\delta_1=\delta_2=0.94\). In the fixed-step
and fixed-path sweeps, the Hybrid LSMC--PDE direct estimator gives smaller
reference errors than plain LSMC in most reported configurations, with the
clearest gains in the low- and moderate-path regimes. These experiments
illustrate the practical value of integrating out the residual one-dimensional
asset uncertainty before performing the volatility-state regression.

The numerical study is deliberately limited to direct estimators and to one fixed
gDMR parameter set. The high-resolution LSMC values are used as references, not
as exact prices. A full study of low-biased estimators, broader parameter
sensitivity, and parameter uncertainty is left for future work.

Future work will therefore focus on extending the analysis in several directions. A natural next step is to study low-biased estimators within the same generalized framework. It would also be of interest to examine alternative contract structures and to study broader parameter regimes under the generalized volatility dynamics. More generally, the present results suggest that the Hybrid LSMC--PDE methodology is a promising computational approach for Bermudan pricing problems in richer multifactor stochastic-volatility settings.



\backmatter

\bmhead{Acknowledgements}
The authors thank Anatoliy Malyarenko for helpful discussions.

\section*{Declarations}

\subsection*{Funding}
Not applicable.

\subsection*{Conflict of interest/Competing interests}
The authors declare that they have no competing interests.

\subsection*{Ethics approval and consent to participate}
Not applicable.

\subsection*{Consent for publication}
Not applicable.

\subsection*{Data availability}
The numerical data used to generate the tables and figures are included in the manuscript source files.

\subsection*{Materials availability}
Not applicable.

\subsection*{Code availability}
Code is available from the corresponding author upon reasonable request.

\subsection*{Author contribution}
Mara Kalicanin Dimitrov and Ying Ni contributed to the conceptualization, mathematical analysis, numerical design, interpretation of results, and manuscript preparation. All authors read and approved the final manuscript.
% Appendix price tables generated from canonical numerical CSV sources.
% Prices are rounded uniformly to three decimal places.
\begin{appendices}

\section{Direct prices for the numerical experiments}\label{app:direct-prices}
This appendix reports the direct Bermudan put prices corresponding to the numerical experiments in Section~\ref{sec:numerical}. All entries are rounded to three decimal places. The label LSMC denotes the plain least-squares Monte Carlo estimator, and Hybrid denotes the Hybrid LSMC--PDE estimator.

\subsection{Baseline numerical specification}

\begin{table}[htbp]
\centering
\caption{Direct reference prices for the baseline numerical specification. Prices are rounded to three decimal places.}
\label{tab:app-reference-prices}
\footnotesize
\setlength{\tabcolsep}{4pt}
\begin{tabularx}{0.86\textwidth}{@{}l *{5}{>{\centering\arraybackslash}X}@{}}
\toprule
Reference setting & $K=70$ & $K=80$ & $K=90$ & $K=100$ & $K=110$ \\
\midrule
\(1200\) steps, \(1{,}200{,}000\) paths & 2.523 & 4.299 & 6.942 & 10.682 & 15.744 \\
\botrule
\end{tabularx}
\end{table}

\begin{table}[htbp]
\centering
\caption{Direct prices in the fixed-path step-sweep experiment with \(20{,}000\) paths. Prices are rounded to three decimal places.}
\label{tab:app-step-sweep-prices-20k}
\footnotesize
\setlength{\tabcolsep}{3pt}
\begin{tabular}{lrrrr}
\toprule
Method/step & 24 & 48 & 72 & 96 \\
\midrule
LSMC \(K=70\) & 2.741 & 2.612 & 2.579 & 2.536 \\
Hybrid \(K=70\) & 2.711 & 2.575 & 2.547 & 2.565 \\
LSMC \(K=80\) & 4.634 & 4.433 & 4.409 & 4.396 \\
Hybrid \(K=80\) & 4.528 & 4.345 & 4.307 & 4.339 \\
LSMC \(K=90\) & 7.341 & 7.129 & 7.068 & 7.080 \\
Hybrid \(K=90\) & 7.172 & 6.957 & 6.905 & 6.956 \\
LSMC \(K=100\) & 11.138 & 10.807 & 10.816 & 10.837 \\
Hybrid \(K=100\) & 10.898 & 10.673 & 10.607 & 10.679 \\
LSMC \(K=110\) & 16.241 & 15.920 & 15.894 & 15.798 \\
Hybrid \(K=110\) & 15.953 & 15.747 & 15.676 & 15.768 \\
\botrule
\end{tabular}
\end{table}

\begin{table}[htbp]
\centering
\caption{Direct prices in the fixed-path step-sweep experiment with \(60{,}000\) paths. Prices are rounded to three decimal places.}
\label{tab:app-step-sweep-prices-60k}
\footnotesize
\setlength{\tabcolsep}{3pt}
\begin{tabular}{lrrrr}
\toprule
Method/step & 24 & 48 & 72 & 96 \\
\midrule
LSMC \(K=70\) & 2.718 & 2.612 & 2.579 & 2.533 \\
Hybrid \(K=70\) & 2.721 & 2.585 & 2.557 & 2.543 \\
LSMC \(K=80\) & 4.531 & 4.434 & 4.362 & 4.318 \\
Hybrid \(K=80\) & 4.547 & 4.364 & 4.326 & 4.307 \\
LSMC \(K=90\) & 7.195 & 7.055 & 7.033 & 6.949 \\
Hybrid \(K=90\) & 7.207 & 6.987 & 6.940 & 6.918 \\
LSMC \(K=100\) & 10.996 & 10.800 & 10.794 & 10.712 \\
Hybrid \(K=100\) & 10.953 & 10.715 & 10.662 & 10.639 \\
LSMC \(K=110\) & 16.050 & 15.788 & 15.853 & 15.806 \\
Hybrid \(K=110\) & 16.026 & 15.803 & 15.749 & 15.725 \\
\botrule
\end{tabular}
\end{table}

\begin{table}[htbp]
\centering
\caption{Direct prices in the fixed 48-step path-sweep experiment. Prices are rounded to three decimal places.}
\label{tab:app-path-sweep-48-prices}
\footnotesize
\setlength{\tabcolsep}{2.5pt}
\begin{tabular}{lrrrrrrr}
\toprule
Method/path count & 250 & 1,000 & 5,000 & 10,000 & 20,000 & 40,000 & 60,000 \\
\midrule
LSMC \(K=70\) & 3.288 & 2.946 & 2.673 & 2.625 & 2.612 & 2.576 & 2.612 \\
Hybrid \(K=70\) & 3.089 & 2.702 & 2.608 & 2.554 & 2.575 & 2.570 & 2.585 \\
LSMC \(K=80\) & 5.425 & 4.795 & 4.545 & 4.452 & 4.433 & 4.369 & 4.434 \\
Hybrid \(K=80\) & 5.079 & 4.533 & 4.427 & 4.310 & 4.345 & 4.357 & 4.364 \\
LSMC \(K=90\) & 8.601 & 7.433 & 7.241 & 7.169 & 7.129 & 7.088 & 7.055 \\
Hybrid \(K=90\) & 7.946 & 7.187 & 7.090 & 6.903 & 6.957 & 6.992 & 6.987 \\
LSMC \(K=100\) & 13.454 & 11.404 & 11.072 & 10.918 & 10.807 & 10.881 & 10.800 \\
Hybrid \(K=100\) & 11.870 & 10.910 & 10.833 & 10.601 & 10.673 & 10.730 & 10.715 \\
LSMC \(K=110\) & 19.278 & 16.843 & 15.985 & 16.057 & 15.920 & 15.921 & 15.788 \\
Hybrid \(K=110\) & 17.120 & 15.962 & 15.889 & 15.660 & 15.747 & 15.819 & 15.803 \\
\botrule
\end{tabular}
\end{table}

\begin{table}[htbp]
\centering
\caption{Direct prices in the fixed 60-step path-sweep experiment. Prices are rounded to three decimal places.}
\label{tab:app-path-sweep-60-prices}
\footnotesize
\setlength{\tabcolsep}{2.5pt}
\begin{tabular}{lrrrrrrr}
\toprule
Method/path count & 250 & 1,000 & 5,000 & 10,000 & 20,000 & 40,000 & 60,000 \\
\midrule
LSMC \(K=70\) & 3.452 & 2.884 & 2.738 & 2.665 & 2.635 & 2.619 & 2.587 \\
Hybrid \(K=70\) & 3.125 & 2.639 & 2.653 & 2.568 & 2.563 & 2.549 & 2.570 \\
LSMC \(K=80\) & 5.950 & 4.762 & 4.629 & 4.553 & 4.484 & 4.387 & 4.411 \\
Hybrid \(K=80\) & 5.034 & 4.396 & 4.470 & 4.321 & 4.318 & 4.321 & 4.342 \\
LSMC \(K=90\) & 8.936 & 7.227 & 7.270 & 7.222 & 7.120 & 7.045 & 7.084 \\
Hybrid \(K=90\) & 7.836 & 6.957 & 7.115 & 6.910 & 6.910 & 6.939 & 6.957 \\
LSMC \(K=100\) & 13.773 & 11.673 & 11.026 & 10.955 & 10.910 & 10.850 & 10.823 \\
Hybrid \(K=100\) & 11.761 & 10.596 & 10.831 & 10.604 & 10.610 & 10.666 & 10.677 \\
LSMC \(K=110\) & 18.962 & 17.240 & 15.920 & 15.991 & 15.918 & 15.952 & 15.860 \\
Hybrid \(K=110\) & 16.953 & 15.621 & 15.863 & 15.659 & 15.680 & 15.755 & 15.757 \\
\botrule
\end{tabular}
\end{table}

\end{appendices}

\bibliography{references}

\end{document}
